\chapter*{\textcolor{gray}{Conclusion}}
L'étude des neurones, de leur membrane et des potentiels électriques qui les régissent s'appuie sur une modélisation précise. Le modèle de Hodgkin-Huxley décrit la dynamique des flux ioniques à l'origine de l'activité neuronale. Nos analyses et simulations ont mis en lumière la capacité de ce modèle à capturer les phénomènes biologiques complexes, offrant un cadre clair pour comprendre le fonctionnement du cerveau.
Le modèle que nous avons étudié, centré sur l'état stationnaire, n'est qu'une partie du travail de Hodgkin et Huxley.\\

Dans leur travail fondamental de 1952, ils ont également étudié la propagation du courant le long de l'axone du neurone, en plus de décrire la dynamique des potentiels d'action au niveau de la membrane. Leur modèle complet inclut aussi la propagation spatiale et temporelle des signaux nerveux, en prenant en compte les variations dynamiques des courants ioniques, comme ceux du sodium et du potassium, pour modéliser la transmission des impulsions nerveuses.\\

Ces approches ouvrent la voie à une compréhension approfondie des réseaux neuronaux et de leurs applications, allant des neurosciences computationnelles aux systèmes intelligents.
\begin{itemize}
	\item \textbf{Modèles réalistes vs simplifiés} : Les modèles multi-compartimentaux et à conductance réaliste offrent une précision accrue pour simuler des neurones complexes, intégrant la géométrie dendritique et les dynamiques ioniques détaillées. À l'opposé, des modèles comme FitzHugh-Nagumo ou Morris-Lecar privilégient la simplicité théorique, facilitant l'analyse mathématique des comportements neuronaux fondamentaux.
	\item \textbf{Wilson-Cowan vs Hodgkin-Huxley} : Le modèle de Wilson-Cowan excelle dans l'étude des dynamiques collectives de populations neuronales, capturant les oscillations corticales et les comportements à grande échelle. Plus abstrait que Hodgkin-Huxley, il se concentre sur les interactions excitateur-inhibiteur au niveau des réseaux, plutôt que sur les mécanismes cellulaires détaillés.
\end{itemize}